\section{Development process}

Development proceeded by turning live failures into invariants. First, early
agents confused observed bargaining field names with outgoing action fields and
allowed multiple clients to race on one turn. Strict schemas, a single play
loop, and offline edge-case tests made interface reliability part of the policy.

Second, V3 remained legal but cycled for 99 rounds under an unlimited horizon
and extreme discount asymmetry. Repeated-state detection and explicit
stall-dependent acceptance repaired termination without accepting known
negative surplus. Third, V11 requested 24 negotiation games through a blocking,
concurrency-four call but returned after 53, losing 44.84 rating points. V33
therefore uses one queue assignment followed by an authoritative audit.

Finally, later batches repeatedly showed role reversals: positive negotiation
totals hid losing buyers, and persuasion buyers often hid losing sellers. V28--
V30 isolated the buyer rollback, Bob recovery, and seller pooling hypotheses.
V33 combined these as fixed role branches and replaced terminal-only deception
with the static bound and deterministic frequency schedule in
Eqs.~\ref{eq:obedience}--\ref{eq:pacing}. These changes were motivated by logged
failure modes, not exhaustive leaderboard parameter search.
